\chapter{Prospective New Zealand Operational Readiness Boundary}\label{app:prospective}

\section{Frozen registration and disabled status}

The accepted prospective registration is version 5, with a ten-year validity horizon from 26 October 2026 to 26 October 2036 (UTC, left-closed and right-open). Its original identity remains pinned. Final target acquisition and scientific scoring are deferred until at least 90 days after the endpoint. The readiness record has \texttt{operations\_scope=readiness\_only}, \texttt{live\_enabled=false} and a null prospective scientific score. No live capture, issue, missing-day mutation, scheduler installation or external alert was performed in that readiness cycle or in this consolidation.

The original registration passed a read-only verification. A separate hard-disabled wrapper supports status/planning and refuses live execution; editing configuration cannot activate the pinned implementation. Its mock runner uses disposable sentinel-marked roots, handcrafted responses and simulated clocks. Those are operational engineering checks, not simulated earthquake futures, operational forecasts or evidence of forecast skill.

\section{Retained mock outcomes and the process-custody gap}

The final author roster passed 22 of 22 disposable checks; an independent repetition passed the same 22. Two earlier failed rosters remain recorded. The checks cover contention, duplicate entry, simulated clock boundaries, held-state recovery, receipt failures, timeouts and specified child-process cases. They certify only their stated disabled/mock-only scope, not production archive custody, power-loss durability, hostile-race security or unattended reliability.

The original independent source review identified a \texttt{CreateProcessW} crash window: a parent can terminate after creating a suspended child but before assigning it to the kill-on-close Job Object. The tested assigned/running-child cleanup does not close that preassignment gap. The later isolated fixed mock uses atomic job membership through a process-creation attribute before \texttt{CreateProcessW}; this targets the historical preassignment gap. It is not production integration. A live implementation still requires separately reviewed archive custody, environment, launcher integration and operational recovery. This monograph does not supply or enable that system.

\subsection{The retained defect-repair successor}

The final operations-repair report retains three execution epochs. Epoch 2 passed 22 of 28 cases, with six failures and 23 terminal exact child handles observed out of 25 expected. Diagnostic epoch 4 passed 28 of 28 and observed all 25 expected handles terminal. Reviewed final epoch 5 again passed 28 of 28 with all 25 terminal and the required trace gates satisfied. Four expected-zero cases are not applicable; a missing expected child is explicitly unknown, not not-applicable. Seven embedded handcrafted controls and 18 independently authored admission controls passed.

The accepted review inspected saved outputs and source; it was not an independent process rerun. Earlier intermittency remains unexplained because the first-run stderr and disposable metadata were discarded. Later success does not prove a timing-cause repair or erase the failures. The successor remains mock-only and live-disabled; no collector, registered archive, scheduler, alert or production integration was enabled. These distinctions are bound to the final probe receipt and independent saved-output review in the claim ledger.

\section{Calendar uncertainty is retained}

At the historical preflight, the first registered day's capture and issue deadlines had elapsed and the local registered archive was absent. That day remains unknown pending independent authoritative reconciliation; absence of one local path is not proof that no external copy existed. No backfill or missing-day command was run. The corrected registration interpretation permits request start and full response completion through the upper capture boundary inclusively. The earlier half-open wording is preserved as superseded history, not adopted as the final rule.

The prospective programme therefore contributes a frozen protocol, disabled implementation and bounded readiness evidence. It contributes zero live forecasts from these cycles and no prospective scientific score. Its operational gate remains separate from all retrospective scientific findings in this volume.
